\chapter[MDP и уравнение Беллмана]{Марковские процессы принятия решений и уравнение Беллмана}
\label{ch:mdp}

\begin{quote}
\itshape
Теория динамического программирования зиждется на единственном наблюдении: будущее независимо от прошлого при условии знания настоящего.
\end{quote}

\bigskip

В главе~\ref{ch:bandits} мы изучали задачу многорукого бандита, в которой действия не имеют последствий помимо немедленного вознаграждения. Реальное принятие решений богаче: усердная учёба сегодня влияет на результат завтрашнего экзамена, который в свою очередь влияет на карьерные перспективы в следующем году. Выборы агента формируют мир, в котором он затем существует. Чтобы строго рассуждать о подобных задачах, нам необходима математическая рамка, улавливающая взаимодействие состояний, действий, вознаграждений и переходов\,---\,рамка, достаточно ёмкая, чтобы кодировать шахматы, локомоцию роботов, диалоговые системы и выравнивание языковых моделей одинаково.

Этой рамкой является \emph{марковский процесс принятия решений} (MDP)\index{Markov Decision Process}. Введённый Ричардом Беллманом в 1950-х годах и формализованный Говардом (1960) и Путерманом (1994), MDP стал стандартной моделью последовательного принятия решений в условиях неопределённости. Практически каждый алгоритм, изучаемый в этой книге\,---\,динамическое программирование, методы Монте-Карло, временное разностное обучение и методы градиента политики\,---\,можно понять как стратегию решения или приближённого решения MDP.

Эта глава строит рамку MDP с нуля. Мы начинаем с марковского свойства, являющегося как математическим удобством, так и глубоким утверждением о структуре мира. Затем формально определяем MDP, вводим политики и функции ценности, выводим знаменитые уравнения Беллмана и изучаем их теоретические свойства. Глава завершается детальным числовым примером, связывающим все концепции воедино.


% ===========================================================
\section{Марковское свойство}
\label{sec:markov-property}

\subsection*{Мотивация}

Представьте робота, перемещающегося в здании. В любой момент у робота есть положение, ориентация и скорость. Чтобы решить, куда двигаться дальше, нужно ли ему помнить каждый шаг с момента включения? Интуитивно — нет\,---\,знания текущего положения и скорости должно быть достаточно для планирования следующего хода. Эта интуиция и есть марковское свойство.

Более формально рассмотрим последовательность случайных величин $S_0, A_0, R_1, S_1, A_1, R_2, \ldots$, представляющих состояние, действие и вознаграждение на каждом шаге. Полная история до момента $t$ — кортеж
\[
H_t = (S_0, A_0, R_1, S_1, A_1, R_2, \ldots, S_{t-1}, A_{t-1}, R_t, S_t).
\]
В принципе следующее состояние $S_{t+1}$ может зависеть от всей этой истории. Марковское свойство утверждает, что важно только текущее состояние $S_t$.

\begin{definition}[Марковское свойство]
\label{def:markov-property}
\index{Markov property}
Случайный процесс $(S_t)_{t \ge 0}$ удовлетворяет \textbf{марковскому свойству}, если для всех $t \ge 0$ и всех состояний $s, s' \in \cS$ и действий $a \in \cA$,
\begin{equation}
\label{eq:markov-property}
\Prob(S_{t+1} = s' \mid S_t = s, A_t = a, S_{t-1}, A_{t-1}, \ldots, S_0, A_0)
= \Prob(S_{t+1} = s' \mid S_t = s, A_t = a).
\end{equation}
Если~\eqref{eq:markov-property} выполняется, говорят, что текущее состояние $S_t$ является \textbf{достаточной статистикой} для истории $H_t$.
\end{definition}

\begin{remark}
\index{Markov property!joint reward-state}
В постановке RL мы дополнительно требуем, чтобы совместное распределение следующего состояния \emph{и} вознаграждения удовлетворяло марковскому свойству:
\[
\Prob(S_{t+1} = s', R_{t+1} = r \mid H_t) = \Prob(S_{t+1} = s', R_{t+1} = r \mid S_t, A_t).
\]
Именно эта форма используется на протяжении всей книги.
\end{remark}

\subsection*{Примеры}

\begin{example}[Модель погоды]
\label{ex:weather}
Предположим, что завтрашняя погода зависит только от сегодняшней, но не от прошлой недели. Если мы кодируем состояние как текущую погоду (солнечно, облачно или дождливо), марковское свойство выполняется и динамика описывается матрицей переходов $3 \times 3$. Добавление большей детализации\,---\,температуры, влажности, давления\,---\,может дать ещё лучшее марковское состояние.
\end{example}

\begin{example}[Навигация робота]
\label{ex:robot-nav}
Следующее положение мобильного робота зависит от текущего положения и поданной команды скорости. Если мы представляем состояние как $(x, y, \theta, v)$---положение, курс, скорость\,---\,динамика примерно марковская. Принципиально, что чисто позиционное состояние $(x, y)$ \emph{не} было бы марковским, поскольку будущее движение зависит также от текущей скорости.
\end{example}

\subsection*{Когда выполняется марковское свойство?}

Для многих естественных представлений состояния марковское свойство является хорошим приближением, а не точной истиной. Погодные паттерны имеют дальнодействующие зависимости; положение робота зависит от износа суставов. На практике мы выбираем представление состояния максимально информативным при условии вычислительной приемлемости.

Когда марковское свойство нарушается\,---\,то есть когда агент может наблюдать лишь \emph{частичный} вид истинного состояния\,---\,мы входим в область \emph{частично наблюдаемых MDP} (POMDP)\index{POMDP}. POMDPs существенно сложнее решать, и мы кратко возвращаемся к ним в главе~\ref{ch:agents}. Пока мы предполагаем, что марковское свойство выполняется точно.

\subsection*{Почему марковское свойство важно?}

Марковское свойство имеет глубокое следствие: любое оптимальное правило принятия решений можно сформулировать как функцию только текущего состояния. Без него оптимальная политика могла бы нуждаться в обращении к постоянно растущей истории. Это сделало бы и анализ, и вычисления неприемлемыми. С ним мы можем записать краткие условия оптимальности\,---\,уравнения Беллмана\,---\,и разработать эффективные алгоритмы для нахождения или приближения оптимального поведения.


% ===========================================================
\section{Формальное определение MDP}
\label{sec:mdp-definition}

\begin{definition}[Марковский процесс принятия решений]
\label{def:mdp}
\index{MDP!definition}
\textbf{(Дисконтированный, конечный) марковский процесс принятия решений} — это кортеж
\begin{equation}
\label{eq:mdp-tuple}
\cM = (\cS,\, \cA,\, p,\, \gamma),
\end{equation}
где:
\begin{description}[leftmargin=!,labelwidth=3.5cm]
  \item[$\cS$] — конечное множество \textbf{состояний}\index{state space};
  \item[$\cA$] — конечное множество \textbf{действий}\index{action space} (пишем $\cA(s)$ для действий, доступных в состоянии $s$);
  \item[$p(s',r \mid s, a)$] — \textbf{ядро перехода-вознаграждения}\index{transition kernel}: вероятность того, что выполнение действия $a$ в состоянии $s$ приводит к следующему состоянию $s'$ и вознаграждению $r \in \cR \subset \R$;
  \item[$\gamma \in [0, 1)$] — \textbf{коэффициент дисконтирования}\index{discount factor}.
\end{description}
Множество вознаграждений $\cR$ считается конечным (или счётным) и ограниченным: существует $R_{\max} > 0$ такое, что $|r| \le R_{\max}$ для всех $r \in \cR$. Ядро $p$ удовлетворяет условию нормировки
\begin{equation}
\label{eq:normalisation}
\sum_{s' \in \cS} \sum_{r \in \cR} p(s', r \mid s, a) = 1
\qquad \forall s \in \cS,\ \forall a \in \cA(s).
\end{equation}
\end{definition}

Функция $p$ — единственный наиболее важный объект в MDP: она кодирует всё, что может делать среда. На каждом шаге $t$ агент наблюдает $S_t = s$, выбирает $A_t = a$, и среда извлекает $(S_{t+1}, R_{t+1}) \sim p(\cdot, \cdot \mid s, a)$.

\subsection*{Производные величины}

Работать напрямую с совместным ядром $p(s', r \mid s, a)$ иногда громоздко. Удобнее три производные величины.

\begin{proposition}[Производные величины]
\label{prop:derived}
\index{transition probability}
\index{expected reward}
Для любых $s \in \cS$, $a \in \cA(s)$, $s' \in \cS$:
\begin{align}
p(s' \mid s, a) &= \sum_{r \in \cR} p(s', r \mid s, a), \label{eq:state-tp} \\[4pt]
r(s, a) &= \E[R_{t+1} \mid S_t = s, A_t = a]
           = \sum_{s' \in \cS} \sum_{r \in \cR} r\, p(s', r \mid s, a), \label{eq:exp-reward-sa} \\[4pt]
r(s, a, s') &= \E[R_{t+1} \mid S_t = s, A_t = a, S_{t+1} = s']
              = \sum_{r \in \cR} r\, \frac{p(s', r \mid s, a)}{p(s' \mid s, a)},
              \quad p(s' \mid s, a) > 0. \label{eq:exp-reward-sas}
\end{align}
\end{proposition}

\begin{proof}
\textbf{Уравнение~\eqref{eq:state-tp}.} Маргинальная вероятность перехода в $s'$ получается суммированием совместного распределения по всем значениям вознаграждения:
\begin{align*}
p(s' \mid s, a)
&= \Prob(S_{t+1} = s' \mid S_t = s, A_t = a)\\
&= \sum_{r \in \cR} \Prob(S_{t+1} = s', R_{t+1} = r \mid S_t = s, A_t = a)
= \sum_{r \in \cR} p(s', r \mid s, a).
\end{align*}

\textbf{Уравнение~\eqref{eq:exp-reward-sa}.} По закону полного математического ожидания,
\[
\E[R_{t+1} \mid S_t = s, A_t = a]
= \sum_{r \in \cR} r\, \Prob(R_{t+1} = r \mid S_t = s, A_t = a).
\]
Маргинализуя по $s'$:
\[
\Prob(R_{t+1} = r \mid S_t = s, A_t = a)
= \sum_{s' \in \cS} p(s', r \mid s, a).
\]
Подстановка даёт~\eqref{eq:exp-reward-sa}.

\textbf{Уравнение~\eqref{eq:exp-reward-sas}.} По определению условного математического ожидания,
\[
\E[R_{t+1} \mid S_t = s, A_t = a, S_{t+1} = s']
= \sum_{r \in \cR} r\, \Prob(R_{t+1} = r \mid S_t = s, A_t = a, S_{t+1} = s').
\]
Применяя теорему Байеса,
\begin{align*}
&\Prob(R_{t+1} = r \mid S_t = s, A_t = a, S_{t+1} = s')\\
&\qquad= \frac{\Prob(S_{t+1} = s', R_{t+1} = r \mid S_t = s, A_t = a)}{\Prob(S_{t+1} = s' \mid S_t = s, A_t = a)}
= \frac{p(s', r \mid s, a)}{p(s' \mid s, a)}.
\end{align*}
Подстановка даёт~\eqref{eq:exp-reward-sas}.
\end{proof}


% ===========================================================
\section{Эпизоды, доходы и коэффициент дисконтирования}
\label{sec:returns}

\subsection*{Эпизодические и непрерывные задачи}

\index{episodic task}\index{continuing task}
Задачи RL делятся на два вида. В \textbf{эпизодической задаче} взаимодействие агента и среды естественно разбивается на конечные отрезки, называемые \emph{эпизодами}: партия в шахматы, захват объекта рукой робота, диалог с пользователем. Каждый эпизод заканчивается в \textbf{терминальном состоянии}\index{terminal state}, после чего среда сбрасывается. В \textbf{непрерывной задаче} взаимодействие продолжается бесконечно: торговый алгоритм, контроллер энергосети, рекомендательная система.

\subsection*{Доход}

Какой бы ни была задача, цель агента — максимизировать суммарное собранное вознаграждение. Формализуем это с помощью \emph{дохода}.

\begin{definition}[Доход]
\label{def:return}
\index{return}
\textbf{Доход} $G_t$ в момент $t$ — это дисконтированная сумма будущих вознаграждений:
\begin{equation}
\label{eq:return}
G_t = \sum_{k=0}^{\infty} \gamma^k R_{t+k+1}
    = R_{t+1} + \gamma R_{t+2} + \gamma^2 R_{t+3} + \cdots.
\end{equation}
Для эпизодической задачи, заканчивающейся на шаге $T$, вознаграждения после $T$ равны нулю, поэтому~\eqref{eq:return} сводится к $G_t = \sum_{k=0}^{T-t-1} \gamma^k R_{t+k+1}$.
\end{definition}

\subsection*{Рекуррентное свойство и его доказательство}

Доход удовлетворяет простой, но принципиальной рекурсии.

\begin{theorem}[Рекуррентный доход]
\label{thm:recursive-return}
\index{return!recursive property}
Для всех $t$,
\begin{equation}
\label{eq:return-recursion}
G_t = R_{t+1} + \gamma G_{t+1}.
\end{equation}
\end{theorem}

\begin{proof}
Начиная с определения~\eqref{eq:return}:
\[
G_t = R_{t+1} + \gamma R_{t+2} + \gamma^2 R_{t+3} + \cdots
    = R_{t+1} + \gamma \bigl(R_{t+2} + \gamma R_{t+3} + \gamma^2 R_{t+4} + \cdots\bigr)
    = R_{t+1} + \gamma G_{t+1}. \qedhere
\]
\end{proof}

Это однострочное доказательство обманчиво важно. Оно гласит, что ценность нахождения в состоянии прямо сейчас равна немедленному вознаграждению плюс дисконтированная ценность следующего состояния\,---\,основа каждого уравнения Беллмана и алгоритма временного разностного обучения.

\subsection*{Ограниченность дисконтированного дохода}

\begin{proposition}[Ограниченный доход]
\label{prop:bounded-return}
Если $|R_t| \le R_{\max}$ для всех $t$ и $\gamma \in [0, 1)$, то
\begin{equation}
\label{eq:bounded-return}
|G_t| \le \frac{R_{\max}}{1 - \gamma}.
\end{equation}
\end{proposition}

\begin{proof}
По неравенству треугольника и формуле суммы геометрического ряда:
\[
|G_t| \le \sum_{k=0}^{\infty} \gamma^k |R_{t+k+1}|
         \le R_{\max} \sum_{k=0}^{\infty} \gamma^k
         = \frac{R_{\max}}{1-\gamma}. \qedhere
\]
\end{proof}

\subsection*{Роль коэффициента дисконтирования}

Параметр $\gamma$ интерполирует между двумя крайностями. При $\gamma = 0$ агент \emph{близоруко}\index{myopic agent}: он заботится только о немедленном вознаграждении $R_{t+1}$. При $\gamma \to 1$ агент \emph{дальновиден}\index{far-sighted agent}: вознаграждения в далёком будущем важны почти так же, как немедленные. На практике $\gamma$ выбирается близко к $1$ (например, $0.99$) для задач, требующих долгосрочного планирования, и ближе к $0$ для задач, где наиболее информативны недавние вознаграждения.

Дисконтирование имеет и финансовую интерпретацию: $\gamma$ — коэффициент процента за шаг, поэтому $\gamma^k$ — приведённая стоимость вознаграждения, полученного через $k$ шагов.

\begin{example}[Вычисление дохода]
\label{ex:return-calc}
\index{return!worked example}
Предположим, агент наблюдает вознаграждения $R_1 = 1,\ R_2 = 2,\ R_3 = -1,\ R_4 = 3$, после чего следует терминальное состояние (все последующие вознаграждения равны нулю), коэффициент дисконтирования $\gamma = 0.9$.

Вычисляя в обратном порядке, начиная с $t = 4$:
\begin{align*}
G_4 &= 0 \quad \mbox{(терминальное)}, \\
G_3 &= R_4 + \gamma G_4 = 3 + 0.9 \times 0 = 3.000, \\
G_2 &= R_3 + \gamma G_3 = -1 + 0.9 \times 3 = -1 + 2.7 = 1.700, \\
G_1 &= R_2 + \gamma G_2 = 2 + 0.9 \times 1.7 = 2 + 1.53 = 3.530, \\
G_0 &= R_1 + \gamma G_1 = 1 + 0.9 \times 3.53 = 1 + 3.177 = 4.177.
\end{align*}
Заметим, что вычисление в обратном порядке через рекурсию~\eqref{eq:return-recursion} эффективно: одного прохода по траектории достаточно.
\end{example}


% ===========================================================
\section{Политики}
\label{sec:policies}

\emph{Политика}\index{policy} определяет поведение агента. Формально она отображает состояния в распределения над действиями.

\begin{definition}[Политика]
\label{def:policy}
\index{policy!stochastic}
\index{policy!deterministic}
\textbf{Стохастическая политика} $\pi$ — это семейство распределений вероятностей над действиями, индексированных состояниями:
\begin{equation}
\label{eq:policy}
\pi(a \mid s) = \Prob(A_t = a \mid S_t = s), \qquad s \in \cS,\ a \in \cA(s).
\end{equation}
\textbf{Детерминированная политика} — частный случай, в котором существует функция $\pi : \cS \to \cA$ такая, что $\pi(a \mid s) = \ind\{a = \pi(s)\}$.
\end{definition}

\begin{remark}
Мы используем один и тот же символ $\pi$ как для стохастического, так и для детерминированного случаев; контекст всегда прояснит, что имеется в виду. Все детерминированные политики являются стохастическими политиками (с точечными распределениями), но не наоборот.
\end{remark}

\begin{example}[Политики в решётке $2 \times 2$]
\label{ex:policy-gridworld}
Рассмотрим четырёхклеточную решётку с клетками $\{(0,0),(0,1),(1,0),(1,1)\}$ и действиями $\{\mbox{Вверх, Вниз, Влево, Вправо}\}$. Детерминированная политика может отображать $(0,0) \mapsto \mbox{Вправо}$, $(0,1) \mapsto \mbox{Вверх}$, $(1,0) \mapsto \mbox{Вправо}$, $(1,1)\mapsto\mbox{Вверх}$. Стохастическая политика может назначать равную вероятность $0.25$ каждому действию в каждом состоянии\,---\,равномерное случайное блуждание.
\end{example}

При фиксированной политике $\pi$ MDP сводится к марковскому процессу с вознаграждением: переходы между состояниями управляются
\[
p_\pi(s' \mid s) = \sum_{a \in \cA(s)} \pi(a \mid s)\, p(s' \mid s, a),
\]
а ожидаемое вознаграждение из состояния $s$ равно $r_\pi(s) = \sum_{a} \pi(a \mid s)\, r(s, a)$.


% ===========================================================
\section{Функции ценности}
\label{sec:value-functions}

Функции ценности измеряют, насколько хорошо находиться в данном состоянии или выполнять данное действие в данном состоянии при следовании политике $\pi$.

\begin{definition}[Функции ценности состояния и действия]
\label{def:value-functions}
\index{value function!state-value}
\index{value function!action-value}
Пусть $\pi$ — политика. \textbf{Функция ценности состояния} $V_\pi : \cS \to \R$ равна
\begin{equation}
\label{eq:state-value}
V_\pi(s) = \E_\pi[G_t \mid S_t = s],
\end{equation}
а \textbf{функция ценности действия} $Q_\pi : \cS \times \cA \to \R$ равна
\begin{equation}
\label{eq:action-value}
Q_\pi(s, a) = \E_\pi[G_t \mid S_t = s,\, A_t = a].
\end{equation}
Здесь $\E_\pi[\,\cdot\,]$ обозначает математическое ожидание по распределению над траекториями, порождаемому $\pi$ и $p$.
\end{definition}

\subsection*{Связь между $V_\pi$ и $Q_\pi$}

\begin{proposition}
\label{prop:v-q-relation}
\index{value function!V-Q relationship}
Для любой политики $\pi$ и состояния $s \in \cS$,
\begin{equation}
\label{eq:v-q-relation}
V_\pi(s) = \sum_{a \in \cA(s)} \pi(a \mid s)\, Q_\pi(s, a).
\end{equation}
\end{proposition}

\begin{proof}
Обусловливая на действие, выбранное в момент $t$, и применяя закон полного математического ожидания:
\begin{align*}
V_\pi(s)
&= \E_\pi[G_t \mid S_t = s] \\
&= \sum_{a \in \cA(s)} \Prob_\pi(A_t = a \mid S_t = s)\,
   \E_\pi[G_t \mid S_t = s, A_t = a] \\
&= \sum_{a \in \cA(s)} \pi(a \mid s)\, Q_\pi(s, a). \qedhere
\end{align*}
\end{proof}

Интуитивно $V_\pi(s)$ — взвешенное по политике среднее ценностей действий: ожидаемый доход агента — это среднее того, что дало бы каждое действие, взвешенное по вероятности, которую политика назначает каждому действию.

\subsection*{Числовой пример: MDP с тремя состояниями}

\begin{example}[Вычисление функции ценности]
\label{ex:value-computation}
\index{value function!worked example}
Рассмотрим MDP на рис.~\ref{fig:three-state-mdp} с тремя состояниями $\{s_1, s_2, s_3\}$, двумя действиями $\{a_1, a_2\}$, доступными только в $s_1$ (состояния $s_2$ и $s_3$ — поглощающие с нулевым вознаграждением), и коэффициентом дисконтирования $\gamma = 0.5$.

\begin{figure}[ht]
\centering
\begin{tikzpicture}[
    >=Stealth,
    node distance=3.0cm,
    state/.style={draw, circle, minimum size=1.0cm, thick},
    absorb/.style={draw, circle, minimum size=1.0cm, thick, double}
  ]
  \node[state]  (s1) {$s_1$};
  \node[absorb, right=3.2cm of s1] (s2) {$s_2$};
  \node[absorb, below=2.2cm of s2] (s3) {$s_3$};

  % действие a1
  \draw[->, bend left=18, thick]
        (s1) to node[above] {$a_1,\ p{=}0.8,\ r{=}5$} (s2);
  \draw[->, bend right=10, thick]
        (s1) to node[right] {$a_1,\ p{=}0.2,\ r{=}1$} (s3);

  % действие a2
  \draw[->, bend right=30, thick, dashed]
        (s1) to node[below left] {$a_2,\ p{=}1,\ r{=}3$} (s3);

  % петли на поглощающих состояниях
  \draw[->, loop right, thick] (s2) to node[right] {$r{=}0$} (s2);
  \draw[->, loop right, thick] (s3) to node[right] {$r{=}0$} (s3);
\end{tikzpicture}
\caption{MDP с тремя состояниями. Действие $a_1$ (сплошная линия) ведёт в $s_2$ с вероятностью $0.8$ и вознаграждением $5$, или в $s_3$ с вероятностью $0.2$ и вознаграждением $1$. Действие $a_2$ (пунктир) детерминированно ведёт в $s_3$ с вознаграждением $3$. Состояния $s_2$ и $s_3$ — поглощающие.}
\label{fig:three-state-mdp}
\end{figure}

Поскольку $s_2$ и $s_3$ поглощающие с нулевым будущим вознаграждением, $V_\pi(s_2) = V_\pi(s_3) = 0$ для любой политики.

Для $s_1$ ценности действий вычисляются непосредственно из определения~\eqref{eq:action-value}:
\begin{align*}
Q_\pi(s_1, a_1)
&= \E[R_{t+1} + \gamma G_{t+1} \mid S_t = s_1, A_t = a_1] \\
&= 0.8 \cdot (5 + 0.5 \cdot 0) + 0.2 \cdot (1 + 0.5 \cdot 0)
 = 4.0 + 0.2 = 4.2, \\
Q_\pi(s_1, a_2)
&= 1.0 \cdot (3 + 0.5 \cdot 0) = 3.0.
\end{align*}

Предположим, политика $\pi(a_1 \mid s_1) = 0.6$, $\pi(a_2 \mid s_1) = 0.4$. Тогда по~\eqref{eq:v-q-relation}:
\[
V_\pi(s_1) = 0.6 \times 4.2 + 0.4 \times 3.0 = 2.52 + 1.20 = 3.72.
\]
При жадной политике ($\pi(a_1 \mid s_1) = 1$) получаем $V_\pi(s_1) = 4.2$, что больше.
\end{example}


% ===========================================================
\section{Уравнение Беллмана}
\label{sec:bellman-equation}

Уравнение Беллмана\index{Bellman equation} — центральное тождество теории RL. Оно выражает функцию ценности как условие самосогласованности: ценность состояния должна равняться ожидаемому немедленному вознаграждению плюс дисконтированная ценность состояния-преемника.

\subsection*{Вывод для $V_\pi$}

\begin{theorem}[Уравнение ожиданий Беллмана для $V_\pi$]
\label{thm:bellman-v}
\index{Bellman equation!for $V_\pi$}
Для любой политики $\pi$ и любого состояния $s \in \cS$,
\begin{equation}
\label{eq:bellman-v}
V_\pi(s)
= \sum_{a \in \cA(s)} \pi(a \mid s)
  \sum_{s' \in \cS} \sum_{r \in \cR}
  p(s', r \mid s, a)\bigl[r + \gamma V_\pi(s')\bigr].
\end{equation}
\end{theorem}

\begin{proof}
Начинаем с определения~\eqref{eq:state-value} и применяем рекуррентный доход~\eqref{eq:return-recursion}:
\[
V_\pi(s)
= \E_\pi[G_t \mid S_t = s]
= \E_\pi[R_{t+1} + \gamma G_{t+1} \mid S_t = s].
\]
По линейности математического ожидания:
\[
V_\pi(s) = \E_\pi[R_{t+1} \mid S_t = s]
           + \gamma\, \E_\pi[G_{t+1} \mid S_t = s].
\]
Применяем свойство башни к второму члену, обусловливая на $S_{t+1}$:
\[
\E_\pi[G_{t+1} \mid S_t = s]
= \E_\pi\!\bigl[\E_\pi[G_{t+1} \mid S_{t+1}] \mid S_t = s\bigr]
= \E_\pi[V_\pi(S_{t+1}) \mid S_t = s].
\]
Следовательно, $V_\pi(s) = \E_\pi[R_{t+1} + \gamma V_\pi(S_{t+1}) \mid S_t = s]$.  Явно раскрывая математическое ожидание по действиям и переходам:
\begin{align*}
V_\pi(s)
&= \sum_{a} \pi(a \mid s)\,
   \E\bigl[R_{t+1} + \gamma V_\pi(S_{t+1}) \mid S_t = s, A_t = a\bigr] \\
&= \sum_{a} \pi(a \mid s)
   \sum_{s'} \sum_{r} p(s', r \mid s, a)\bigl[r + \gamma V_\pi(s')\bigr]. \qedhere
\end{align*}
\end{proof}

\subsection*{Уравнение Беллмана для $Q_\pi$}

\begin{proposition}[Уравнение Беллмана для $Q_\pi$]
\label{prop:bellman-q}
\index{Bellman equation!for $Q_\pi$}
\begin{equation}
\label{eq:bellman-q}
Q_\pi(s, a)
= \sum_{s' \in \cS} \sum_{r \in \cR} p(s', r \mid s, a)
  \Bigl[r + \gamma \sum_{a' \in \cA(s')} \pi(a' \mid s')\, Q_\pi(s', a')\Bigr].
\end{equation}
\end{proposition}

\begin{proof}
Начиная с $Q_\pi(s,a) = \E_\pi[R_{t+1} + \gamma G_{t+1} \mid S_t=s, A_t=a]$ и применяя свойство башни, обусловленное на $(S_{t+1}, R_{t+1})$, затем используя $\E_\pi[G_{t+1} \mid S_{t+1} = s'] = V_\pi(s') = \sum_{a'} \pi(a' \mid s') Q_\pi(s', a')$.
\end{proof}

\subsection*{Резервная диаграмма}

Уравнение Беллмана имеет естественное графическое представление, известное как \emph{резервная диаграмма}\index{backup diagram}, показанная на рис.~\ref{fig:backup-v}. Пустые кружки представляют состояния; закрашенные кружки — пары состояние-действие. Корень — состояние $s$; диаграмма разветвляется к действиям по $\pi$, затем далее к состояниям-преемникам по $p$.

\begin{figure}[ht]
\centering
\begin{tikzpicture}[
    >=Stealth,
    level 1/.style={sibling distance=3.2cm, level distance=1.6cm},
    level 2/.style={sibling distance=1.6cm, level distance=1.6cm},
    state/.style={draw, circle, minimum size=0.65cm, fill=white, thick},
    sa/.style={draw, circle, minimum size=0.55cm, fill=black!80, thick},
  ]
  % корневое состояние
  \node[state] (root) {$s$}
    child {
      node[sa] (a1) {}
      child { node[state] {$s'$} edge from parent node[left,  font=\footnotesize] {$p,r$} }
      child { node[state] {$s''$} edge from parent node[right, font=\footnotesize] {$p,r$} }
      edge from parent node[left, font=\footnotesize] {$\pi$}
    }
    child {
      node[sa] (a2) {}
      child { node[state] {$s'''$} edge from parent node[left, font=\footnotesize] {$p,r$} }
      child { node[state] {$\cdots$} edge from parent node[right, font=\footnotesize] {$p,r$} }
      edge from parent node[right, font=\footnotesize] {$\pi$}
    };
\end{tikzpicture}
\caption{Резервная диаграмма для уравнения Беллмана для $V_\pi(s)$. Пустые кружки — состояния; закрашенные кружки — пары состояние-действие. Рёбра с меткой $\pi$ показывают выбор действия; рёбра с меткой $p,r$ — переходы среды с соответствующим вознаграждением.}
\label{fig:backup-v}
\end{figure}

\subsection*{Матричная форма для конечных MDP}

Когда $\cS$ конечно и $n = |\cS|$ состояний, уравнение~\eqref{eq:bellman-v} можно записать в матрично-векторной форме. Определим:
\begin{itemize}[leftmargin=*]
  \item $\mathbf{v}_\pi \in \R^n$ с $[\mathbf{v}_\pi]_s = V_\pi(s)$;
  \item $\mathbf{r}_\pi \in \R^n$ с $[\mathbf{r}_\pi]_s = r_\pi(s) = \sum_a \pi(a \mid s) r(s,a)$;
  \item $\mathbf{P}_\pi \in \R^{n \times n}$ с $[\mathbf{P}_\pi]_{s,s'} = \sum_a \pi(a \mid s) p(s' \mid s, a)$.
\end{itemize}

Тогда уравнение Беллмана принимает вид линейной системы
\begin{equation}
\label{eq:bellman-matrix}
\mathbf{v}_\pi = \mathbf{r}_\pi + \gamma \mathbf{P}_\pi \mathbf{v}_\pi,
\end{equation}
с единственным решением
\begin{equation}
\label{eq:bellman-matrix-solution}
\mathbf{v}_\pi = (I - \gamma \mathbf{P}_\pi)^{-1} \mathbf{r}_\pi.
\end{equation}
Матрица $(I - \gamma \mathbf{P}_\pi)$ обратима, поскольку $\mathbf{P}_\pi$ — стохастическая матрица (спектральный радиус $1$), поэтому все собственные значения $\gamma \mathbf{P}_\pi$ имеют модуль не более $\gamma < 1$, что делает все собственные значения $(I - \gamma \mathbf{P}_\pi)$ отделёнными от нуля.

\begin{example}[Решение системы Беллмана]
\label{ex:bellman-system}
\index{Bellman equation!worked example}
Вернёмся к примеру~\ref{ex:value-computation}, но теперь предположим, что оба состояния $s_2$ и $s_3$ возвращаются обратно в $s_1$ с вознаграждением $0$, и $s_1$ использует политику $\pi(a_1|s_1)=0.6$, $\pi(a_2|s_1)=0.4$, $\gamma=0.5$.

Матрица переходов под $\pi$ из $s_1$:
\[
p_\pi(s_2 \mid s_1) = 0.6 \times 0.8 = 0.48,\quad
p_\pi(s_3 \mid s_1) = 0.6 \times 0.2 + 0.4 \times 1 = 0.52.
\]
Ожидаемое вознаграждение: $r_\pi(s_1) = 0.6 \times 4.2 + 0.4 \times 3.0 = 3.72$ (как и прежде). Для простоты полагаем $V(s_2) = V(s_3) = 0$ (поглощающие). Уравнение Беллмана для $s_1$:
\[
V_\pi(s_1) = 3.72 + 0.5 (0.48 \cdot 0 + 0.52 \cdot 0) = 3.72.
\]
Это совпадает с прямым вычислением в примере~\ref{ex:value-computation}.
\end{example}


% ===========================================================
\section{Оптимальные функции ценности и уравнение оптимальности Беллмана}
\label{sec:bellman-optimality}

\subsection*{Оптимальные функции ценности}

\begin{definition}[Оптимальные функции ценности]
\label{def:optimal-value}
\index{optimal value function}
\index{optimal policy}
\textbf{Оптимальная функция ценности состояния}:
\begin{equation}
\label{eq:v-star}
V_*(s) = \sup_\pi V_\pi(s), \qquad s \in \cS,
\end{equation}
а \textbf{оптимальная функция ценности действия}:
\begin{equation}
\label{eq:q-star}
Q_*(s, a) = \sup_\pi Q_\pi(s, a), \qquad s \in \cS,\ a \in \cA(s).
\end{equation}
Политика $\pi_*$ является \textbf{оптимальной}, если $V_{\pi_*}(s) = V_*(s)$ для всех $s \in \cS$.
\end{definition}

\subsection*{Упорядочение политик и существование оптимальной политики}

\begin{definition}[Упорядочение политик]
\label{def:policy-ordering}
\index{policy ordering}
Пишем $\pi' \succeq \pi$ и говорим, что $\pi'$ \textbf{не хуже} $\pi$, если
\[
V_{\pi'}(s) \ge V_\pi(s) \qquad \forall s \in \cS.
\]
\end{definition}

\begin{theorem}[Существование оптимальной политики]
\label{thm:optimal-policy-exists}
\index{optimal policy!existence}
Для любого конечного MDP существует по меньшей мере одна детерминированная оптимальная политика $\pi_*$ такая, что $V_{\pi_*}(s) = V_*(s)$ для всех $s \in \cS$.
\end{theorem}

\begin{proof}[Набросок доказательства]
Ключевое наблюдение: политика оптимальна тогда и только тогда, когда в каждом состоянии $s$ она действует жадно относительно $V_*$:
\[
\pi_*(a \mid s) > 0 \implies a \in \argmax_{a' \in \cA(s)} Q_*(s, a').
\]
Такая детерминированная политика всегда существует (разрешая ничьи произвольно). Из уравнений оптимальности Беллмана, выведенных ниже, следует, что $V_{\pi_*} = V_*$.
\end{proof}

Оптимальные функции ценности удовлетворяют собственным уравнениям Беллмана, получаемым заменой взвешенной по политике суммы по действиям на максимум.

\begin{theorem}[Уравнения оптимальности Беллмана]
\label{thm:bellman-optimality}
\index{Bellman optimality equation}
\begin{align}
V_*(s) &= \max_{a \in \cA(s)}
           \sum_{s' \in \cS} \sum_{r \in \cR}
           p(s', r \mid s, a)\bigl[r + \gamma V_*(s')\bigr],
           \label{eq:bellman-opt-v} \\[4pt]
Q_*(s, a) &= \sum_{s' \in \cS} \sum_{r \in \cR}
             p(s', r \mid s, a)
             \Bigl[r + \gamma \max_{a' \in \cA(s')} Q_*(s', a')\Bigr].
             \label{eq:bellman-opt-q}
\end{align}
\end{theorem}

\begin{proof}
Для~\eqref{eq:bellman-opt-v}: поскольку $V_*(s) = \max_\pi V_\pi(s)$, уравнение Беллмана~\eqref{eq:bellman-v} для оптимальной политики $\pi_*$ даёт
\[
V_*(s) = \sum_a \pi_*(a|s) \sum_{s',r} p(s',r|s,a)[r + \gamma V_*(s')].
\]
Поскольку $\pi_*$ жадна относительно $V_*$, взвешенная сумма по $a$ равна максимуму по $a$, что даёт~\eqref{eq:bellman-opt-v}. Уравнение~\eqref{eq:bellman-opt-q} следует аналогично из $Q$-версии уравнения Беллмана.
\end{proof}

\begin{remark}
Уравнения оптимальности Беллмана \emph{нелинейны} относительно $V_*$ (из-за оператора $\max$) и не могут быть решены простым обращением матрицы. Необходимо динамическое программирование и связанные алгоритмы; см. главу~\ref{ch:dp}.
\end{remark}


% ===========================================================
\section{Оператор Беллмана}
\label{sec:bellman-operator}

Уравнения Беллмана можно изучать как задачи нахождения неподвижных точек линейных операторов — перспектива, объясняющая, почему алгоритмы их решения сходятся.

\begin{definition}[Операторы Беллмана]
\label{def:bellman-operators}
\index{Bellman operator}
Пусть $\cB(\cS)$ обозначает банахово пространство ограниченных функций $V : \cS \to \R$, оснащённое супремум-нормой $\|V\|_\infty = \max_{s \in \cS} |V(s)|$.

\textbf{Оператор Беллмана для политики} $\cT_\pi : \cB(\cS) \to \cB(\cS)$:
\begin{equation}
\label{eq:bellman-policy-op}
(\cT_\pi V)(s)
= \sum_{a \in \cA(s)} \pi(a \mid s)
  \sum_{s', r} p(s', r \mid s, a)\bigl[r + \gamma V(s')\bigr].
\end{equation}
\textbf{Оператор оптимальности Беллмана} $\cT_* : \cB(\cS) \to \cB(\cS)$:
\begin{equation}
\label{eq:bellman-opt-op}
(\cT_* V)(s)
= \max_{a \in \cA(s)}
  \sum_{s', r} p(s', r \mid s, a)\bigl[r + \gamma V(s')\bigr].
\end{equation}
\end{definition}

В этих обозначениях уравнения Беллмана читаются как $V_\pi = \cT_\pi V_\pi$ и $V_* = \cT_* V_*$: функции ценности являются \emph{неподвижными точками} соответствующих операторов Беллмана.

\begin{theorem}[$\gamma$-сжатие]
\label{thm:contraction}
\index{Bellman operator!contraction}
Оба оператора $\cT_\pi$ и $\cT_*$ являются $\gamma$-сжатиями в супремум-норме: для любых $V, W \in \cB(\cS)$,
\begin{align}
\|\cT_\pi V - \cT_\pi W\|_\infty &\le \gamma \|V - W\|_\infty, \label{eq:contraction-pi} \\
\|\cT_* V - \cT_* W\|_\infty &\le \gamma \|V - W\|_\infty. \label{eq:contraction-star}
\end{align}
\end{theorem}

\begin{proof}
Доказываем~\eqref{eq:contraction-pi}; доказательство~\eqref{eq:contraction-star} аналогично. Фиксируем $s \in \cS$:
\begin{align*}
|(\cT_\pi V)(s) - (\cT_\pi W)(s)|
&= \Bigl|\sum_a \pi(a|s) \sum_{s',r} p(s',r|s,a)\,\gamma\bigl(V(s') - W(s')\bigr)\Bigr| \\
&\le \sum_a \pi(a|s) \sum_{s',r} p(s',r|s,a)\,\gamma\, |V(s') - W(s')| \\
&\le \gamma\, \|V - W\|_\infty
     \underbrace{\sum_a \pi(a|s) \sum_{s',r} p(s',r|s,a)}_{=\,1}.
\end{align*}
Беря максимум по $s$, получаем~\eqref{eq:contraction-pi}.

Для~\eqref{eq:contraction-star} используем неравенство $|\max_a f(a) - \max_a g(a)| \le \max_a |f(a) - g(a)|$ и действуем аналогично.
\end{proof}

\begin{corollary}[Теорема Банаха о неподвижной точке]
\label{cor:banach}
\index{Banach fixed-point theorem}
Поскольку $\cB(\cS)$ — полное метрическое пространство (конечномерное) и операторы $\cT_\pi$, $\cT_*$ являются сжатиями ($\gamma < 1$), теорема Банаха о неподвижной точке гарантирует:
\begin{enumerate}
  \item Каждый оператор имеет \emph{единственную} неподвижную точку: $V_\pi$ и $V_*$ соответственно.
  \item Итерирование оператора, начиная с произвольного $V^{(0)} \in \cB(\cS)$:
        \[V^{(k+1)} = \cT V^{(k)}\]
        сходится к неподвижной точке с геометрической скоростью:
        \[
        \|V^{(k)} - V_\pi\|_\infty \le \gamma^k \|V^{(0)} - V_\pi\|_\infty \to 0.
        \]
\end{enumerate}
\end{corollary}

Это следствие является теоретической основой \emph{оценки политики} в динамическом программировании (глава~\ref{ch:dp}) и обучения TD(0) (глава~\ref{ch:td}): оба можно понять как итеративное применение (приближения) оператора Беллмана.


% ===========================================================
\section{Числовой пример: MDP студента}
\label{sec:student-mdp}

Теперь разберём полный пример, вдохновлённый MDP студента из Саттона и Барто (2018).\index{student MDP}

\begin{example}[MDP студента]
\label{ex:student-mdp}
Студент может находиться в одном из четырёх состояний:
\begin{itemize}[leftmargin=*]
  \item $s_1$: \textbf{Занятия} — посещение лекции.
  \item $s_2$: \textbf{Facebook} — просмотр социальных сетей.
  \item $s_3$: \textbf{Учёба} — подготовка к экзамену.
  \item $s_4$: \textbf{Сон} — сон (терминальное/поглощающее состояние).
\end{itemize}
Действия и переходная динамика приведены в таблице~\ref{tab:student-mdp}, а MDP изображён на рис.~\ref{fig:student-mdp}. Используем коэффициент дисконтирования $\gamma = 0.9$.

\begin{table}[ht]
\centering
\caption{Переходная динамика MDP студента. Каждая строка задаёт действие в текущем состоянии, следующее состояние, вероятность и полученное вознаграждение.}
\label{tab:student-mdp}
\begin{tabular}{lllrr}
\toprule
Состояние & Действие & Следующее состояние & Вер.\ $p$ & Вознаграждение $r$ \\
\midrule
$s_1$ (Занятия) & посещать & $s_3$ (Учёба) & 1.0 & $-2$ \\
$s_1$ (Занятия) & отвлечься & $s_2$ (Facebook) & 1.0 & $-1$ \\
$s_2$ (Facebook) & выйти & $s_1$ (Занятия) & 1.0 & $-1$ \\
$s_2$ (Facebook) & продолжить & $s_2$ (Facebook) & 1.0 & $-1$ \\
$s_3$ (Учёба) & сдать & $s_4$ (Сон) & 1.0 & $+10$ \\
$s_3$ (Учёба) & учиться ещё & $s_3$ (Учёба) & 1.0 & $-2$ \\
\bottomrule
\end{tabular}
\end{table}

\begin{figure}[ht]
\centering
\begin{tikzpicture}[
    >=Stealth, thick,
    state/.style={draw, rounded corners, minimum width=2.0cm, minimum height=0.8cm},
    every edge/.style={->, thick}
  ]
  \node[state] (c)  at (0, 0)    {Занятия $s_1$};
  \node[state] (fb) at (-3, -2)  {Facebook $s_2$};
  \node[state] (st) at (3, -2)   {Учёба $s_3$};
  \node[state, double] (sl) at (3, -4.5) {Сон $s_4$};

  \draw (c) edge[bend left=15] node[above right, font=\small] {посещать, $-2$} (st);
  \draw (c) edge[bend right=15] node[above left, font=\small] {отвлечься, $-1$} (fb);
  \draw (fb) edge[bend right=15] node[below right, font=\small] {выйти, $-1$} (c);
  \draw (fb) edge[loop left] node[left, font=\small] {продолжить, $-1$} (fb);
  \draw (st) edge node[right, font=\small] {сдать, $+10$} (sl);
  \draw (st) edge[loop right] node[right, font=\small] {учиться ещё, $-2$} (st);
\end{tikzpicture}
\caption{MDP студента. Рёбра помечены действием и немедленным вознаграждением. Терминальное состояние Сон ($s_4$, двойная граница) является поглощающим.}
\label{fig:student-mdp}
\end{figure}

\paragraph{Вычисление функций ценности при оптимальной политике.}
Оптимальная стратегия очевидно — посещать занятия и затем сдать экзамен: посещать $s_1$, достичь $s_3$, немедленно сдать в $s_4$.

При этой политике $\pi_*$:
\begin{align*}
V_*(s_4) &= 0 \quad \mbox{(поглощающее)}, \\
V_*(s_3) &= r(\mbox{сдать}) + \gamma V_*(s_4) = 10 + 0.9 \times 0 = 10.0, \\
V_*(s_1) &= r(\mbox{посещать}) + \gamma V_*(s_3) = -2 + 0.9 \times 10 = -2 + 9 = 7.0.
\end{align*}
Таким образом, оптимальная ценность состояния «Занятия» равна $7.0$.

\paragraph{Вычисление ценностей действий.}
В $s_1$:
\begin{align*}
Q_*(s_1, \mbox{посещать})   &= -2 + 0.9 \times 10 = 7.0, \\
Q_*(s_1, \mbox{отвлечься}) &= -1 + 0.9 \times V_*(s_2).
\end{align*}
Для $V_*(s_2)$ заметим, что из Facebook оптимальное действие — \emph{выйти} (возвращаясь к Занятиям, имеющим ценность $7.0$), поэтому:
\[
V_*(s_2) = -1 + 0.9 \times 7.0 = -1 + 6.3 = 5.3.
\]
Следовательно:
\[
Q_*(s_1, \mbox{отвлечься}) = -1 + 0.9 \times 5.3 = -1 + 4.77 = 3.77.
\]
Поскольку $Q_*(s_1, \mbox{посещать}) = 7.0 > 3.77 = Q_*(s_1, \mbox{отвлечься})$, жадная политика действительно выбирает «посещать», подтверждая наши прежние рассуждения.

\paragraph{Проверка уравнения оптимальности Беллмана в $s_2$.}
Уравнение оптимальности Беллмана~\eqref{eq:bellman-opt-v} гласит:
\[
V_*(s_2)
= \max\bigl\{Q_*(s_2, \mbox{выйти}),\, Q_*(s_2, \mbox{продолжить})\bigr\}.
\]
Имеем:
\begin{align*}
Q_*(s_2, \mbox{выйти}) &= -1 + 0.9 \times 7.0 = 5.3, \\
Q_*(s_2, \mbox{продолжить}) &= -1 + 0.9 \times V_*(s_2) = -1 + 0.9 \times 5.3 = 3.77.
\end{align*}
Действительно $\max\{5.3, 3.77\} = 5.3 = V_*(s_2)$. \checkmark
\end{example}


% ===========================================================
\section{Резюме}
\label{sec:mdp-summary}

\begin{itemize}[leftmargin=*]
  \item \textbf{Марковское свойство} утверждает, что будущее независимо от прошлого при условии знания текущего состояния. Оно является основой, делающей RL разрешимым.
  \item \textbf{Марковский процесс принятия решений} $(\cS, \cA, p, \gamma)$ — стандартная рамка для последовательного принятия решений. Ядро $p(s', r \mid s, a)$ кодирует всю динамику.
  \item \textbf{Доход} $G_t = \sum_k \gamma^k R_{t+k+1}$ фиксирует дисконтированное суммарное будущее вознаграждение. Его рекуррентное свойство $G_t = R_{t+1} + \gamma G_{t+1}$ является зерном уравнений Беллмана.
  \item \textbf{Функции ценности} $V_\pi(s)$ и $Q_\pi(s,a)$ измеряют ожидаемый доход под политикой $\pi$. Они связаны соотношением $V_\pi(s) = \sum_a \pi(a|s) Q_\pi(s,a)$.
  \item \textbf{Уравнения ожиданий Беллмана} выражают функции ценности как самосогласованные неподвижные точки. Для малых MDP они могут быть решены точно методами линейной алгебры.
  \item \textbf{Уравнения оптимальности Беллмана} характеризуют $V_*$ и $Q_*$ через оператор $\max$. Их решения непосредственно дают оптимальную политику.
  \item \textbf{Операторы Беллмана} $\cT_\pi$ и $\cT_*$ являются $\gamma$-сжатиями в супремум-норме. По теореме Банаха о неподвижной точке итерирование любого из операторов из произвольной начальной точки геометрически сходится к единственной неподвижной точке. Это теоретический фундамент динамического программирования и обучения TD.
\end{itemize}


% ===========================================================
\section*{Упражнения}
\addcontentsline{toc}{section}{Упражнения}

\begin{exercise}
\label{ex:markov-state}
Приведите пример задачи, в которой необработанное сенсорное наблюдение (например, один пиксель) \emph{не} является марковским состоянием, а дополненное состояние (например, последние $k$ кадров) является. Объясните, почему дополненное состояние удовлетворяет марковскому свойству.
\end{exercise}

\begin{exercise}
\label{ex:return-constant}
Рассмотрим непрерывную задачу, в которой агент всегда получает вознаграждение $c > 0$ на каждом шаге и $\gamma \in [0,1)$.
\begin{enumerate}[label=(\alph*)]
  \item Вычислите $G_t$ в замкнутой форме.
  \item Покажите, что $G_t = c/(1-\gamma)$.
  \item Объясните, почему условие $\gamma < 1$ необходимо для конечности дохода.
  \item Что происходит при $\gamma = 1$?
\end{enumerate}
\end{exercise}

\begin{exercise}
\label{ex:q-from-p}
Используя определения из предложения~\ref{prop:derived}, проверьте, что
\[
r(s, a) = \sum_{s' \in \cS} p(s' \mid s, a)\, r(s, a, s').
\]
Это показывает, что трёхаргументное ожидаемое вознаграждение $r(s,a,s')$ более информативно, чем $r(s,a)$.
\end{exercise}

\begin{exercise}
\label{ex:action-values-comp}
В MDP студента (пример~\ref{ex:student-mdp}), предположим, студент следует политике, при которой в состоянии $s_3$ он учится ещё с вероятностью $0.5$ и сдаёт с вероятностью $0.5$.  Используйте $\gamma = 0.9$.
\begin{enumerate}[label=(\alph*)]
  \item Запишите уравнение Беллмана для $V_\pi(s_3)$ под этой смешанной политикой.
  \item Решите относительно $V_\pi(s_3)$.
  \item Сравните с $V_*(s_3) = 10$ и объясните разницу.
\end{enumerate}
\end{exercise}

\begin{exercise}
\label{ex:bellman-q-proof}
Докажите уравнение Беллмана для $Q_\pi$ (предложение~\ref{prop:bellman-q}) подробно, начиная с определения $Q_\pi(s,a) = \E_\pi[G_t \mid S_t = s, A_t = a]$ и рекуррентного дохода~\eqref{eq:return-recursion}.
\end{exercise}

\begin{exercise}
\label{ex:matrix-bellman}
Рассмотрим MDP с двумя состояниями $\cS = \{s_1, s_2\}$, одним действием, матрицей переходов
\[
P_\pi = \begin{pmatrix} 0.3 & 0.7 \\ 0.6 & 0.4 \end{pmatrix},
\]
вектором ожидаемых вознаграждений $\mathbf{r}_\pi = (2, 5)^\top$ и $\gamma = 0.8$.
\begin{enumerate}[label=(\alph*)]
  \item Запишите матричное уравнение Беллмана~\eqref{eq:bellman-matrix}.
  \item Явно вычислите $(I - \gamma P_\pi)$.
  \item Решите относительно $\mathbf{v}_\pi$ по формуле~\eqref{eq:bellman-matrix-solution}.
\end{enumerate}
\end{exercise}

\begin{exercise}
\label{ex:contraction-star}
Завершите доказательство теоремы~\ref{thm:contraction} для оператора оптимальности $\cT_*$. \emph{Подсказка}: используйте тождество $|\max_a f(a) - \max_a g(a)| \le \max_a |f(a) - g(a)|$.
\end{exercise}

\begin{exercise}
\label{ex:convergence-rate}
Начиная с $V^{(0)} \equiv 0$ и итерируя $\cT_*$ с $\gamma = 0.9$ и $R_{\max} = 1$:
\begin{enumerate}[label=(\alph*)]
  \item Дайте верхнюю оценку $\|V^{(k)} - V_*\|_\infty$ через $k$ и $\gamma$.
  \item Сколько итераций необходимо, чтобы гарантировать $\|V^{(k)} - V_*\|_\infty \le 0.01$?
\end{enumerate}
\end{exercise}

\begin{exercise}
\label{ex:optimal-policy-unique}
Приведите пример конечного MDP, имеющего бесконечно много оптимальных политик, но единственную оптимальную функцию ценности $V_*$. Убедитесь, что все оптимальные политики достигают $V_*$.
\end{exercise}

\begin{exercise}
\label{ex:pomdp-teaser}
Рассмотрим решётку $1 \times 3$ (клетки $1$, $2$, $3$). Агент может двигаться влево или вправо. Клетка $3$ — цель (вознаграждение $+1$, терминальная). Агент \emph{не} наблюдает текущую клетку; он получает лишь сигнал о том, столкнулся ли он со стеной.
\begin{enumerate}[label=(\alph*)]
  \item Объясните, почему наблюдение не является марковским состоянием.
  \item Определите дополненное состояние (используя историю наблюдений), восстанавливающее марковское свойство.
  \item Опишите оптимальную политику в терминах этого дополненного состояния.
\end{enumerate}
\end{exercise}
